\chapter{相位问题与实倒空间迭代}
\section{相位缺失与不可区分变换}
完整结构因子可写为 $F_h=A_he^{\iu\phi_h}$。实验在通常条件下给出与 $A_h^2$ 成比例的观测，$\phi_h$ 未被直接测量。周期密度的反演式为
\begin{equation}
\rho(\vct x)=\frac1V\sum_{\vct h}F_{\vct h}e^{-2\pi\iu\vct h^{\mathsf T}\vct x}.
\label{eq:rho}
\end{equation}
没有相位就不能直接相加重建密度。给定同一组幅度，任取一组相位通常都产生某个实空间函数；只有极少数满足原子性、组成和对称性。

整体平移 $\vct x_j\mapsto\vct x_j+\vct o$ 使
\begin{equation}
F_h\mapsto e^{2\pi\iu\vct h^{\mathsf T}\vct o}F_h,
\end{equation}
但不改变 $|F_h|^2$。无异常散射时，反演把 $F_h$ 变成 $F_h^*$，同样不改变强度。因此，合成结构恢复误差必须在等价原点、反演和相同元素原子置换下比较。直接比较两组坐标表会把同一个解误判为失败。

\section{Patterson 函数与重原子信息}
用 $|F_h|^2$ 替代式~\eqref{eq:rho} 的复结构因子，得到 Patterson 函数
\begin{equation}
P(\vct u)=\frac1V\sum_h|F_h|^2e^{-2\pi\iu\vct h^{\mathsf T}\vct u}.
\end{equation}
这里 $\vct u$ 为分数位移。对实电子密度、忽略异常散射的情况，利用傅里叶卷积定理可证明
\begin{equation}
P(\vct u)=\int_{\mathrm{cell}}\rho_{\mathrm{cart}}(\vct r)
\rho_{\mathrm{cart}}(\vct r+\mat A\vct u)\dif^3\vct r.
\end{equation}
若原子密度足够集中，非原点峰对应原子间矢量 $\vct r_j-\vct r_i$，峰权重近似与两原子的散射能力乘积相关。含金属的结构因此常能先显现重原子之间的矢量关系。Patterson 的优点是无需相位；其困难是 $N$ 个原子产生约 $N(N-1)$ 个非零矢量，峰会重叠，且矢量集合不总能唯一决定结构。

\section{直接法的基本代数关系}
为说明相位关系的来源，先考虑由相同、彼此不重叠的原子密度组成的结构。原子密度平方仍集中在相同中心，可近似写成
\begin{equation}
\rho^2(\vct r)\simeq C\,\widetilde\rho(\vct r),
\end{equation}
其中 $\widetilde\rho$ 具有相同的原子中心但不同的原子形状因子。对两边作傅里叶变换，乘积转为卷积，得到 Sayre 类型关系~\cite{sayre1952}
\begin{equation}
F_h\simeq\theta_h\sum_k F_kF_{h-k}.
\label{eq:sayre}
\end{equation}
对以强反射为主的项，右端复数向量之和的方向倾向于 $\phi_h$。令权重 $w_k$ 包含幅度与统计可靠度，切线更新形式为
\begin{equation}
\phi_h\leftarrow\operatorname{atan2}\left[
\sum_kw_k\sin(\phi_k+\phi_{h-k}),\;
\sum_kw_k\cos(\phi_k+\phi_{h-k})\right].
\label{eq:tangent}
\end{equation}
这一关系不把单条强反射变成独立的相位观测，而是把若干相位联结成三相位不变量。$\phi_h-\phi_k-\phi_{h-k}$ 中的整体原点相位相互抵消~\cite{direct_methods}。

对元素差异很大的结构，简单的等原子关系需要修改。现代求解方法可结合 Patterson 信息、归一化幅度、峰筛选和多起点搜索。配套源码没有实现完整的传统直接法关系网络；式~\eqref{eq:sayre} 与式~\eqref{eq:tangent} 在这里提供相位约束的数学基础。

\section{两个空间中的约束集合}
设 $\mathcal F$ 为离散傅里叶变换，观测幅度给出倒易空间集合
\begin{equation}
\mathcal M=\{F:\;|F_h|=A_h,\;h\in\Omega\},
\end{equation}
其中 $\Omega$ 是观测反射集合。对应的幅度投影为
\begin{equation}
[P_{\mathcal M}(F)]_h=
\begin{cases}
A_hF_h/|F_h|,&h\in\Omega,\ |F_h|>0,\\
A_h,&h\in\Omega,\ |F_h|=0,\\
F_h,&h\notin\Omega.
\end{cases}
\end{equation}
实空间集合可以表达非负性、原子稀疏性或已知片段。直接交替投影
\begin{equation}
\rho^{(t+1)}=P_{\mathcal A}\mathcal F^{-1}P_{\mathcal M}\mathcal F\rho^{(t)}
\end{equation}
面对的是非凸集合，一般不能保证收敛到真实结构，也可能停留在错误的局部一致解。因此，多起点、独立评分和后续精修都是求解过程的一部分。

\section{电荷翻转迭代}
电荷翻转对低于阈值的密度改变符号~\cite{chargeflip2004}：
\begin{equation}
g^{(t)}(\vct x)=
\begin{cases}
\rho^{(t)}(\vct x),&\rho^{(t)}(\vct x)\ge\delta,\\
-\rho^{(t)}(\vct x),&\rho^{(t)}(\vct x)<\delta.
\end{cases}
\label{eq:chargeflip}
\end{equation}
随后变换到倒易空间，替换幅度，保留新的相位，再回到实空间。这里 $\delta=c\,\operatorname{sd}(\rho)$，$c$ 为输入中记录的阈值系数。原生实现可先定义谱包络 $e_h$，以 $\widetilde A_h=A_h/e_h$ 作为迭代目标；$e_h$ 由已知组成的均方散射能力归一化并设置下限得到。以下步骤与评分中的目标幅度均指 $\widetilde A_h$，关闭锐化时它才等于 $A_h$。每次试验按以下顺序执行：
\begin{enumerate}
\item 用固定随机种子生成实密度，得到满足 Hermitian 关系的初始复系数。
\item 用工作集观测幅度替换已测系数的模；分辨率球外的系数置零。
\item 按式~\eqref{eq:chargeflip} 修改密度，计算新复系数 $G_h$。
\item 在替换幅度之前计算偏差，再实施幅度投影。未测量但位于分辨率球内的系数作为潜变量更新。
\item 保存本次试验中偏差最小的相位；对已测系数乘回 $e_h$，提取周期性局部极大值。未测潜系数仍在迭代尺度上，这一中间图只用于定位峰，不作为已校准的物理密度输出。
\end{enumerate}
投影前的评分为
\begin{equation}
R_{\mathrm{proj}}=\frac{\sum_{h\in\Omega}\big||G_h|-\widetilde A_h\big|}
{\sum_{h\in\Omega}\widetilde A_h}.
\end{equation}
评分必须在幅度投影之前计算，否则投影后的 $|F_h|=\widetilde A_h$ 会使评分恒为零。$R_{\mathrm{proj}}$ 用于比较迭代状态，不是结构正确概率，也不能替代原子模型的精修残差。

\section{网格采样与数值符号}
取网格点 $\vct x_{\vct n}=(n_1/N_1,n_2/N_2,n_3/N_3)$。NumPy 的 \code{fftn} 使用负号指数，因此本报告约定下有
\begin{equation}
\rho_{\vct n}=\frac{1}{V}\operatorname{fftn}(F)_{\vct n},\qquad
F_{\vct h}=V\operatorname{ifftn}(\rho)_{\vct h}.
\end{equation}
若相位迭代使用统一的无量纲密度，常数 $V$ 可相互抵消，但导出的物理密度必须恢复 $1/V$。

为避免正负 Nyquist 指数混叠，各方向必须满足
\begin{equation}
N_i>2\max_{h\in\Omega}|h_i|.
\end{equation}
更大的网格能提高峰位置的插值精度，但不增加实验分辨率。源码的网格大小同时受最大体素数约束，避免 agent 将异常指数展开为不受控制的内存分配。

\section{从峰到原子模型}
原生峰搜索使用周期边界的局部极大值，按晶胞真实度量排除过近峰，再用邻点抛物线作亚网格插值。已知元素列表按原子序数从大到小分配给按峰高排序的位置。这是一种候选构造启发式，不是对元素种类的完整自动判定。

化学先验可帮助筛选候选。例如，给定金属和配体构成时，可检查配位距离、键连通性和单胞组成。但一个几何上合理的片段也可能在错误的相位解中出现。合理顺序是先由衍射提出密度与位置候选，再用明确的化学约束比较候选，最后回到观测强度检验。

\section{空间群的相位检验与原点搜索}
已知一组 $P1$ 相位后，可检验候选空间群是否满足式~\eqref{eq:symF}。若当前密度相对标准设定平移了 $\vct o$，则操作的有效平移为
\begin{equation}
\vct t'=\vct t+(\mat I-\mat R)\vct o.
\end{equation}
候选一致性可用相位残差或复系数差衡量，例如
\begin{equation}
S_G(\vct o)=\sum_{g,h}w_{gh}
\left|F_{\mat R_g^{\mathsf T}h}
-e^{-2\pi\iu h^{\mathsf T}[\vct t_g+(\mat I-\mat R_g)\vct o]}F_h\right|^2.
\end{equation}
这给出了原点搜索的目标函数，同时解释了为什么强度消光相同的空间群仍可由相位区分。完整求解还需要等价原点枚举、伪对称检验、缺失反射处理及候选原子模型的重新循环。原生版本提供相位恢复和消光诊断，但尚未实现这一完整联立空间群搜索；需要相应行为时使用明确选择的官方 SHELXT 路径。
